At a finite public-chance cut, partial counterfactual regret minimization (CFR)
can evaluate successive outcome batches at learner states changed by earlier
batches.  Balanced coverage then lacks the conditional unbiasedness used by
Monte Carlo CFR analyses.  We introduce \emph{target debit}, the one-sided
cumulative underestimation needed to transfer regret on delivered feedback to
the dense target game.  Our central theorem bounds weighted exploitability by
delivered-row regret, target debit, and a native alternating-update correction.
This yields Debit-Budgeted CFR (DB-CFR), which selects a public plan from
completed history, executes it once, and settles its debit afterward.  Complete
coverage has zero debit and supplies exact recovery.  With at most $\ell$
unsettled blocks, hidden debit is bounded by the sum of the $\ell$ largest
predictable block envelopes; for fixed $\ell$ and polynomial output weights,
the resulting $O(T^{-1})$ term preserves the additive-RM/RM+ rate.  A finite
converse gives equal-metadata plans with charges $0$ and $11/8$, proving that a
nontrivial strict-budget decision requires semantic information or bounded
post-execution settlement.  In two released heads-up no-limit hold'em turn
endgames, all 42 registered shallow comparisons favor partial coverage before
complete coverage overtakes it between 32 and 64 rounds.  In public Leduc, a
lagged rule fixed on four development orders beats the best
matched-recurrence-work fixed width on all ten unseen orders (median NashConv
gain $0.0456$) without future candidate replay.
